\usepackage[spanish]{babel}
\usepackage[letterpaper, margin=1in]{geometry}
\usepackage{graphicx}
\graphicspath{{images/}}
\usepackage{times}
\usepackage{setspace}
\usepackage{tabularx}
\usepackage{array}
\usepackage{float}
\usepackage{etoolbox}
\usepackage[colorlinks=true, linkcolor=black, citecolor=black, urlcolor=black]{hyperref}
\usepackage{apacite}
\bibliographystyle{apacite}
\AtBeginDocument{
  \renewcommand{\BOthers}{et al}
  \renewcommand{\BOthersPeriod}{et al.}
}
\usepackage{lipsum} % For placeholder text

% Interlineado y Espacio entre párrafos 
\linespread{1.5}
\setlength{\parskip}{5pt}

% Formato para los capítulos
\usepackage{titlesec}
\titlespacing{\chapter}{0pt}{0pt}{0pt}

% Formato para evitar la separación de palabras 
\tolerance=1
\emergencystretch=\maxdimen
\hyphenpenalty=10000
\hbadness=10000

% Formato de números de páginas
\usepackage{fancyhdr}
\fancyhf{}
\setlength{\headheight}{15pt}

\fancypagestyle{plain}{
  \fancyhf{}
  \fancyfoot[C]{\thepage}
  \renewcommand{\headrulewidth}{0pt}
  \renewcommand{\footrulewidth}{0pt}
}

\renewcommand{\headrulewidth}{0pt}
\renewcommand{\thechapter}{\Roman{chapter}}

% Formato de sangrías y listas
\usepackage{enumitem}
\renewcommand{\labelitemi}{$\bullet$}
\setlength{\parindent}{0.4in}
\setlist[itemize,1]{leftmargin=0.55in}
\setlist[itemize]{leftmargin=0.2in}

% Formato de la tabla de contenido
\usepackage{tocloft}
\addto\captionsspanish{
  \renewcommand{\contentsname}{Índice de Contenido}
  \renewcommand{\listfigurename}{Índice de Figuras}
  \renewcommand{\listtablename}{Índice de Tablas}
}
\renewcommand{\cfttoctitlefont}{\hfill\normalsize\bfseries}
\renewcommand{\cftaftertoctitle}{\hfill}
\renewcommand{\cftloftitlefont}{\hfill\normalsize\bfseries}
\renewcommand{\cftafterloftitle}{\hfill}
\renewcommand{\cftlottitlefont}{\hfill\normalsize\bfseries}
\renewcommand{\cftafterlottitle}{\hfill}
\setlength{\cftbeforetoctitleskip}{-20pt}
\setlength{\cftbeforeloftitleskip}{-20pt}
\setlength{\cftbeforelottitleskip}{-20pt}
\setcounter{secnumdepth}{4}
\setcounter{tocdepth}{4}

% TOC Chapter Style
\cftpagenumbersoff{chapter}
\setlength{\cftbeforechapskip}{5pt}

% TOC Sections Style
\setlength{\cftsecindent}{0pt}
\setlength{\cftsubsecindent}{1.5em}
\setlength{\cftsubsubsecindent}{3em}
\setlength{\cftparaindent}{4.5em}

% Formato de los diferentes niveles de secciones
\newcommand{\upper}[1]{\uppercase{#1}}
\usepackage{titlesec}
\titleformat{\chapter}
{\normalfont\fontsize{12}{15}\bfseries\centering}{}{0.3em}{}[\thispagestyle{empty}]
\titleformat{\section}
{\normalfont\fontsize{12}{15}\bfseries}{\thesection}{1em}{}
\titleformat{\subsection}
{\normalfont\fontsize{12}{15}\bfseries\narrower}{\thesubsection}{1em}{}
\titleformat{\subsubsection}
{\normalfont\fontsize{12}{15}\bfseries\itshape\narrower}{\thesubsubsection}{1em}{}
\titleformat{\paragraph}
{\normalfont\fontsize{12}{15}\itshape\narrower}{\theparagraph}{1em}{}

% Comandos para crear secciones no numeradas con entrada en el índice
\newcommand{\uchapter}[1]{
\refstepcounter{chapter}
\phantomsection
\chapter*{Capítulo \thechapter. {#1}}
\addcontentsline{toc}{chapter}{Capítulo \thechapter. {#1}}
}
\newcommand{\uextra}[2]{%
  \refstepcounter{chapter}%
  \phantomsection
  \chapter*{\textbf{#1 \thechapter. #2}}%
  \addcontentsline{toc}{chapter}{\protect\textbf{#1 \thechapter. #2}}%
}
\newcommand{\usection}[1]{
  \phantomsection
  \section*{#1}
  \addcontentsline{toc}{section}{#1}
}
\newcommand{\usubsection}[1]{
  \phantomsection
  \subsection*{#1.}
  \addcontentsline{toc}{subsection}{#1}
}
\newcommand{\usubsubsection}[1]{
  \phantomsection
  \subsubsection*{#1.}
  \addcontentsline{toc}{subsubsection}{#1}
}
\newcommand{\uparagraph}[1]{
  \phantomsection
  \paragraph*{#1.}
  \addcontentsline{toc}{paragraph}{#1}
}

% Helper para insertar la línea de fuente en captions de figuras.
% Uso: \figsource{<título>}{<autor>}{<año>}{<página>}{<derechos y/o nota opcional>}
% Ejemplo: \figsource{Título del artículo}{Sánchez-Vera}{2019}{30}{Derechos reservados.}
\newcommand{\figsource}[5]{Tomado de \emph{#1} (#2, #3, p.~#4). #5}

% Variante que sigue el formato exacto solicitado:
% "Tomado de TÍTULO, (p. N), por AUTOR, AÑO."
% Uso: \figtakenfrom{<título>}{<página>}{<autor(es)>}{<año>}{<nota opcional>}
\newcommand{\figtakenfrom}[5]{%
  Tomado de \emph{#1}%
  \if\relax\detokenize{#2}\relax\else, (p.~#2)\fi, por #3, #4. #5%
}

% Formato de figuras, tablas y ecuaciones
\usepackage{caption}
\usepackage{chngcntr}

% === Comportamiento de floats (tablas/figuras) ===
% Evita que floats relativamente pequeños se vayan a una página exclusiva.
\renewcommand{\topfraction}{0.9}
\renewcommand{\bottomfraction}{0.8}
\renewcommand{\textfraction}{0.07}
\renewcommand{\floatpagefraction}{0.8}
\setcounter{topnumber}{2}
\setcounter{bottomnumber}{1}
\setcounter{totalnumber}{3}

% === Columnas para tablas con ajuste de texto ===
% Útil con tabularx: columnas tipo X que envuelven texto.
\newcolumntype{Y}{>{\raggedright\arraybackslash}X}
\newcolumntype{L}[1]{>{\raggedright\arraybackslash}p{#1}}
% Especificación por defecto para tablas de 3 columnas:
% 1ra col con ancho fijo, 2da/3ra col se ajustan y envuelven.
\newcommand{\ThreeColWrapSpec}{L{0.28\textwidth} Y Y}

% Espaciado vertical entre filas en tablas (1.0 = normal).
% Sube a 1.15, 1.2, 1.3, etc. según necesidad.
\newcommand{\TableRowStretch}{1.0}

% === Estilo de captions (configurable) ===
% Separador: "Tabla 1." en una línea y el caption en la siguiente.
% Cambia este comando a: period, colon, newline, etc. según necesites.
\newcommand{\TableCaptionLabelSep}{capnewline}
\DeclareCaptionLabelSeparator{capnewline}{\newline}

% Alineación del caption de tablas (izquierda).
% Valores típicos: raggedright, centering, justified
\newcommand{\TableCaptionJustification}{raggedright}

% Tamaño de letra para captions y contenido de tablas.
% Usa normalsize para igualar el texto base del documento.
\newcommand{\TableBodyFontSize}{\small}
\AtBeginEnvironment{tabular}{\TableBodyFontSize}
\AtBeginEnvironment{tabular*}{\TableBodyFontSize}
\AtBeginEnvironment{tabularx}{\TableBodyFontSize}

% Figuras
% Poner el título de la figura en la parte inferior (por defecto), con la palabra
% "Figura" seguida del número en cursiva y el texto del caption también en
% cursiva. Además se fuerza el tamaño de letra del caption para que coincida
% con el del texto normal (no \small ni otro tamaño reducido).
\captionsetup[figure]{
  labelformat=default,
  labelsep=period,
  name=Figura,
  labelfont={it},
  textfont={it,normalfont},
  font=small
}
\renewcommand{\thefigure}{\arabic{figure}}
\counterwithout{figure}{chapter}

% Tablas
\captionsetup[table]{
  labelformat=default,
  labelsep=\TableCaptionLabelSep,
  name=Tabla,
  textfont={it},
  font=small,
  justification=\TableCaptionJustification,
  singlelinecheck=false
}
\renewcommand{\thetable}{\arabic{table}}
\counterwithout{table}{chapter}

% Ecuaciones
\renewcommand{\theequation}{\arabic{equation}}
\counterwithout{equation}{chapter}

\usepackage{calc}
\renewcommand{\cftfigpresnum}{Figura~}
\setlength{\cftfignumwidth}{\widthof{Figura 0. }}
% Queda solo un punto: el que añade `caption` con `labelsep=period`.
% Evitamos añadir otro punto en la lista de figuras.
\renewcommand{\cftfigaftersnum}{}

\renewcommand{\cfttabpresnum}{Tabla~}
\setlength{\cfttabnumwidth}{\widthof{Tabla Ap-Z99 }}
% Evitar punto extra en la lista de tablas, el punto se gestiona con `caption`.
\renewcommand{\cfttabaftersnum}{}
